%Master File:lectures.tex

\lesson{Introduction}
\vspace{-2cm}
\begin{center}
 \includegraphics[height=10cm]{ch01/figure_01_03b}
\end{center}

\keywords{Evolution, DNA, Transciption, Translation}
%%%%%%%%%%%%%%%%%%%%%%% Next Slide %%%%%%%%%%%%%%%%%%%%%%%
\renewcommand{\Outline}{%
\begin{slide}
\section[1]{Outline}

\begin{minipage}{12cm}
  \begin{enumerate}\squeeze
    \outlineitem{Big Data}{bd}
    \outlineitem{Beginnings}{beginnigs}
    \outlineitem{DNA}{dna}
    \outlineitem{Protein Production}{pp}
   \end{enumerate}
\end{minipage}\hfill
\begin{minipage}{10cm}
  \vfill
  \includegraphics[width=10cm]{ch01/figure_01_01}
  \vfill
\end{minipage}
\end{slide}
\addtocounter{outlineitem}{1}
}

\setcounter{outlineitem}{1}
\Outline
\toptarget{firstoutline}
%%%%%%%%%%%%%%%%%%%%%%% Next Slide %%%%%%%%%%%%%%%%%%%%%%%

\begin{slide}
\section{Biology is a Big Data Problem!}

\begin{PauseHighLight}
  \begin{itemize}
  \item Biology is big
    \begin{itemize}
    \item Genomes involve $10^5$--$10^{12}$ nucleotide bases\pause
    \item Number of different proteins can be $10^5$\pause
    \item Number of different cell types $10^2$\pause
    \item Number of species $10^7$\pause
    \end{itemize}
  \item Understanding biology involves a lot of information\pause
  \item This information is stored in big databases\pause
  \end{itemize}
\end{PauseHighLight}

\end{slide}

%%%%%%%%%%%%%%%%%%%%%%% Next Slide %%%%%%%%%%%%%%%%%%%%%%%

\begin{slide}
\section{Interactions}

\begin{PauseHighLight}
  \begin{itemize}
  \item Interactions occur across many levels
    \begin{itemize}
    \item Protein-protein interaction are common place\pause
    \item Protein-DNA interaction underlie a lot of regulation\pause
    \item Protein-environment interaction are how cell sense the world\pause
    \item Species interact in ecosystems\pause
      \begin{itemize}
      \item You are an ecosystem\pauseb
      \item Half the cells in your body aren't yours\pauseb
      \end{itemize}
    \end{itemize}
  \item If $P\approx10^5$ is the number of proteins then the number of possible
    interacting pairs are $P\approx10^{10}$\pause
  \end{itemize}
\end{PauseHighLight}

\end{slide}

%%%%%%%%%%%%%%%%%%%%%%% Next Slide %%%%%%%%%%%%%%%%%%%%%%%

\begin{slide}
\section{Biology Self-Organise}

\begin{PauseHighLight}
  \begin{itemize}
  \item Protein complexes form machines\pause
  \item Life organises into cells\pause
  \item Cells can organise into higher organisms\pause
  \item Organisms form populations that act together\pause
  \item Back to ecosystems\pause
  \end{itemize}
\end{PauseHighLight}

\end{slide}

%%%%%%%%%%%%%%%%%%%%%%% Next Slide %%%%%%%%%%%%%%%%%%%%%%%

\begin{slide}
\section{Life is Dynamic}

\begin{PauseHighLight}
  \begin{itemize}
  \item Cells are chemical machines\pause
  \item They work through complex auto-regulation\pause
  \item Cells divide\pause
  \item Sometimes cell fuse sexually\pause
  \item Life evolves\pause
  \item Almost all aspects of life are governed by control mechanisms\pauseb
  \end{itemize}
\end{PauseHighLight}

\end{slide}

%%%%%%%%%%%%%%%%%%%%%%% Next Slide %%%%%%%%%%%%%%%%%%%%%%%

\begin{slide}
\section{Computational Biology}

\begin{PauseHighLight}
  \begin{itemize}
  \item Making sense of biology is a big data problem\pause
  \item Tools of computer science are essential to
    \begin{itemize}
    \item store the data
    \item make inferences about the data\pause
    \end{itemize}
  \item The tools of abstraction applies to biology as much as to
    computer science\pause
  \item Biology is both infinitely complex\pause{} and yet in some
    cases amazing simple\pauseb
  \end{itemize}
\end{PauseHighLight}


\end{slide}

%%%%%%%%%%%%%%%%%%%%%%% Next Slide %%%%%%%%%%%%%%%%%%%%%%%
\Outline
%%%%%%%%%%%%%%%%%%%%%%% Next Slide %%%%%%%%%%%%%%%%%%%%%%%

\begin{slide}
\section[-1]{Evolution}

\begin{minipage}{0.3\linewidth}
  \tiny{Charles Darwin}\\
  \includegraphics[width=0.9\linewidth]{darwin}\\
  \tiny{Alfred Russel Wallace}\\
  \includegraphics[width=0.9\linewidth]{wallace}
\end{minipage}\hfill
\begin{minipage}{0.65\linewidth}
  \begin{PauseHighLight}
  \begin{itemize}
  \item "Nothing in Biology Makes Sense Except in the Light of
    Evolution"\\
    \hfill--Theodosius Dobzhansky\pause
  \item At every level (from protein to species) we can make sense of
    biology in term of natural selection\pause
  \item Understanding the details of evolution is an on-going story\pause
  \end{itemize}
\end{PauseHighLight}

\end{minipage}
\end{slide}

%%%%%%%%%%%%%%%%%%%%%%% Next Slide %%%%%%%%%%%%%%%%%%%%%%%

\begin{slide}
\section{The Tree of Life}
\begin{minipage}{0.3\linewidth}
  \includegraphics[width=0.9\linewidth]{treeOfLife}\\
\end{minipage}\hfill
\begin{minipage}{0.65\linewidth}
  \begin{PauseHighLight}
  \begin{itemize}
  \item All observations are consistent with life on earth being related
    through a tree like structure\pause
  \item At the molecular level different species share a huge amount in
    common\pause
  \item Biologists use this to transfer knowledge between species\pause
  \end{itemize}
\end{PauseHighLight}

\end{minipage}


\end{slide}

%%%%%%%%%%%%%%%%%%%%%%% Next Slide %%%%%%%%%%%%%%%%%%%%%%%

\begin{slide}
\section{Inheritance}

\begin{PauseHighLight}
  \begin{itemize}
  \item At the time of Darwin (1859) inheritance was \emph{not}
    generally understood\pause
  \item This caused a lot of confusion because blending inheritance
    would seem to remove the variation needed for natural
    selection\pause
  \item Modern inheritance theory was initially proposed by Gregor
    Mendel around 1860, but was not widely known\pause
  \item Around 1900 it was finally widely understood that inheritance
    occurred through the exchange of discrete units
    (genes)\pause \end{itemize}
\end{PauseHighLight}

\end{slide}

%%%%%%%%%%%%%%%%%%%%%%% Next Slide %%%%%%%%%%%%%%%%%%%%%%%

\begin{slide}
\section{Gregor Mendel's Experiment}

\begin{center}
  \includegraphics[width=0.9\linewidth]{mendel}
\end{center}
\end{slide}


%%%%%%%%%%%%%%%%%%%%%%% Next Slide %%%%%%%%%%%%%%%%%%%%%%%

\begin{slide}
  \section{Sexual Reproduction}
  \pb
  \pause
  \pauselevel{=1}
  \begin{center}
    \multipdf[width=\linewidth]{sexualRep}\pause
  \end{center}

\end{slide}

%%%%%%%%%%%%%%%%%%%%%%% Next Slide %%%%%%%%%%%%%%%%%%%%%%%

\begin{slide}
\section[-2]{Mendel's Peas}

\pb
\begin{center}
  \includegraphics[width=25cm,height=15cm]{inheritance}\pauselevel{=1 :5}\pause
  \multido{\ia=0+1,\ib=1+1}{3}{%
  \llap{\includegraphics[width=25cm,height=15cm]{MendelBoxes-\ia}}\mypl{\ib}}
\end{center}
\end{slide}

%%%%%%%%%%%%%%%%%%%%%%% Next Slide %%%%%%%%%%%%%%%%%%%%%%%

\begin{slide}
\section[-1]{Inheritance}

\begin{PauseHighLight}
  \begin{itemize}
  \item Mendel's experiments could be explained if
    \begin{itemize}
    \item traits were carried by discrete units (genes)\pause
    \item each individual carries two copies of a gene (diploidy)\pause
    \item we inherit one copy from each parent\pause
    \item one gene is dominant\pause
    \end{itemize}
  \item These are not always true\pause
  \item Very clever explanation and experiments\pause{}, although the
    results were rigged\pauseb
  \item The mechanism for inheritance remained confused until\ldots\pauseb
  \end{itemize}
\end{PauseHighLight}

\end{slide}



%%%%%%%%%%%%%%%%%%%%%%% Next Slide %%%%%%%%%%%%%%%%%%%%%%%
\Outline
%%%%%%%%%%%%%%%%%%%%%%% Next Slide %%%%%%%%%%%%%%%%%%%%%%%

\begin{slide}
\section[1]{Crick and Watson}

\begin{minipage}{0.5\linewidth}
  \begin{center}
    Discovery of the Double Helix\\
    \includegraphics[width=\linewidth]{DNANature53}
  \end{center}
\end{minipage}\hfill
\begin{minipage}{0.45\linewidth}
  \begin{center}
    \includegraphics[width=\linewidth]{crickWatson}
  \end{center}
\end{minipage}
\end{slide}

%%%%%%%%%%%%%%%%%%%%%%% Next Slide %%%%%%%%%%%%%%%%%%%%%%%

\begin{slide}
\section{Base Pairs}
\begin{minipage}{0.5\linewidth}
  \includegraphics[width=\linewidth]{BasePairing}\\
\end{minipage}\hfill
\begin{minipage}{0.45\linewidth}
  \begin{PauseHighLight}
  \begin{itemize}
  \item Double helix held together by 4 types of bases\pause
  \item Each base, A, C, G, T form \emph{base pairs} A--T or C--G\pause
  \item Amazing digital code\pause
  \end{itemize}
\end{PauseHighLight}

\end{minipage}

\end{slide}



%%%%%%%%%%%%%%%%%%%%%%% Next Slide %%%%%%%%%%%%%%%%%%%%%%%

\begin{slide}
\section{The Genetic Code}

\begin{minipage}{0.2\linewidth}
  \begin{center}
    \includegraphics[width=0.9\linewidth]{dna-cw}\\
  \end{center}
\end{minipage}\hfill
\begin{minipage}{0.75\linewidth}
\begin{quote}
  ``It has not escaped our notice that the specific pairing we have
  postulated immediately suggests a possible copying mechanism for the
  genetic material.''
\end{quote}
\begin{flushright}
    -- Crick and Watson (1953)
\end{flushright}
\end{minipage}
\end{slide}

%%%%%%%%%%%%%%%%%%%%%%% Next Slide %%%%%%%%%%%%%%%%%%%%%%%

\begin{slide}
\section{DNA Replication}

\begin{center}
  \includegraphics[width=0.9\linewidth]{DNA_Replication}
\end{center}
\end{slide}


%%%%%%%%%%%%%%%%%%%%%%% Next Slide %%%%%%%%%%%%%%%%%%%%%%%
\Outline
%%%%%%%%%%%%%%%%%%%%%%% Next Slide %%%%%%%%%%%%%%%%%%%%%%%

\begin{slide}
\section{What Does DNA Do?}

\begin{PauseHighLight}
  \begin{itemize}
  \item Part of DNA codes for Proteins\pause
    \begin{itemize}
    \item Gets copied to messenger RNA through transcription\pause
    \item mRNA then gets translated to proteins\pause
    \end{itemize}
  \item A lot of DNA does \emph{not} code for protein
    \begin{itemize}
    \item Can code for RNA molecules (ribozymes)\pause
    \item Used for controlling the cell\pause
    \item Can be junk---caused by replication\pause
    \item We're not quite sure\pauseb
    \end{itemize}
  \end{itemize}
\end{PauseHighLight}

\end{slide}


%%%%%%%%%%%%%%%%%%%%%%% Next Slide %%%%%%%%%%%%%%%%%%%%%%%

\begin{slide}
\section[-2]{RNA and DNA}

\begin{center}
  \includegraphics[width=0.8\linewidth]{RNAandDNA}
\end{center}

\end{slide}

%%%%%%%%%%%%%%%%%%%%%%% Next Slide %%%%%%%%%%%%%%%%%%%%%%%

\begin{slide}
\section[-1]{Transcription by RNA Polymerase in Action}

\begin{center}
  \includegraphics[width=\linewidth]{rnaPolymerase}
\end{center}

\end{slide}

%%%%%%%%%%%%%%%%%%%%%%% Next Slide %%%%%%%%%%%%%%%%%%%%%%%

\begin{slide}
\section[-2]{Exons and Introns}

\begin{center}
  \includegraphics[width=\linewidth]{exintrons}
\end{center}
\end{slide}


%%%%%%%%%%%%%%%%%%%%%%% Next Slide %%%%%%%%%%%%%%%%%%%%%%%

\begin{slide}
\section[-2]{mRNA to Protein Translation}

\begin{center}
  \includegraphics[width=0.8\linewidth]{ribosomeWorking}
\end{center}
\end{slide}


%%%%%%%%%%%%%%%%%%%%%%% Next Slide %%%%%%%%%%%%%%%%%%%%%%%

\begin{slide}
\section[-2]{RNA Triplet to Amino Acids}
\begin{center}
  \includegraphics[width=0.6\linewidth]{ch01/figure_01_04}
\end{center}
\end{slide}

%%%%%%%%%%%%%%%%%%%%%%% Next Slide %%%%%%%%%%%%%%%%%%%%%%%

\begin{slide}
\section[-2]{DNA to Protein}

\begin{center}
  \includegraphics[width=0.9\linewidth]{DNAtoProtein}
\end{center}
\end{slide}

%%%%%%%%%%%%%%%%%%%%%%% Next Slide %%%%%%%%%%%%%%%%%%%%%%%

\begin{slide}
\section{What Do Proteins Do?}

\begin{PauseHighLight}
  \begin{itemize}
  \item Proteins are the main building block of cells\pause
  \item The can form the skeleton (e.g. lipid layer of membranes)\pause
  \item They acts as chemical machines (enzymes)\pause
  \item They can even act as mechanical machines\pause
  \item They are a large part of the control mechanism\pause
  \item Not everything in the cell is a protein
    \begin{itemize}
      \item E.g. there is DNA and RNA\pause
      \item RNA can fold up like proteins and act as catalysts
        (ribosymes)\pause
    \end{itemize}
  \end{itemize}
\end{PauseHighLight}

\end{slide}

%%%%%%%%%%%%%%%%%%%%%%% Next Slide %%%%%%%%%%%%%%%%%%%%%%%

\begin{slide}
\section{Computational Challenges}

\begin{PauseHighLight}
  \begin{itemize}
  \item A lot of computational challenges had to be addressed to
    understand the genome\pause
  \item Need to find coding regions (exons)\pause
  \item Need to reconstruct DNA sequence from sequencing data\pause
  \item Need to compare different sequences\pause
  \item We will cover these issues during the course\pause{} (and more
    beside)\pauseb
  \end{itemize}
\end{PauseHighLight}

\end{slide}

%%%%%%%%%%%%%%%%%%%%%%% Next Slide %%%%%%%%%%%%%%%%%%%%%%%

\begin{slide}
\section{Conclusion}

\begin{PauseHighLight}
  \begin{itemize}
  \item We've just started to look at molecular biology\pause
  \item We've covered inheritance through DNA\pause
  \item How to produce proteins from DNA\pause
  \item We've still got to look at how cells work\pause
  \end{itemize}
\end{PauseHighLight}

\end{slide}



%%% Local Variables:
%%% TeX-master: "lectures"
%%% End:
